\documentclass[10pt,a4paper]{amsart}
\usepackage[margin=19mm]{geometry}
\usepackage{amsmath,amssymb,mathtools}
\usepackage[scr=rsfs,cal=euler]{mathalfa}
\usepackage{xfrac}
\usepackage[unicode,hidelinks,bookmarks=false]{hyperref}
\newcommand{\ie}{i.e.\ }
\newcommand{\cf}{cf.\ }
\newcommand{\resp}{respectively\ }
\newcommand{\Subsection}{\subsection}
\providecommand{\nobreakdash}{\nobreak\hbox{-}}
\setlength{\emergencystretch}{2em}
\multlinegap0pt
\usepackage{xcolor}
\usepackage[shortlabels]{enumitem}
\usepackage{float}
\newtheorem{assumption}{Assumption}
\newtheorem{theorem}{Theorem}
\def\R{\mathbb{R}} 
\newcommand*\der{\mathop{}\!\mathrm{d}}
\title{Treasure Search Optimization}
\author{A. Sharma}
\date{Source: arXiv:2607.16863v1 (CC BY 4.0)}
\begin{document}
\maketitle

\noindent\emph{Source and licence.} Treasure Search Optimization. Version arXiv:2607.16863v1 (CC BY 4.0), distributed by arXiv under the Creative Commons Attribution 4.0 International licence, http://creativecommons.org/licenses/by/4.0/, as recorded in the arXiv licence metadata for this exact version and shown on the abstract page https://arxiv.org/abs/2607.16863v1. Full licence notice: \texttt{licenses/arxiv-2607.16863-CC-BY-4.0.txt}.\par\smallskip

\noindent\textit{Editorial scope.} Counted unit: the article's theorem thm:steady\_state\_laplace (source0.tex lines 1141--1155). The article's complete original proof of it is delivered with it (source0.tex lines 1156--1441, one \texttt{proof} environment of 286 lines). No mathematics is written, repaired or re-derived: the statement and the proof are the article's own lines, in the article's own notation, and no retained line is substituted or re-flowed.  Retained uncounted as the source-local context the proof needs: lines 1022--1024; lines 1133--1138.  Nothing was omitted from any retained span, so the package carries no omission record.  No retained line contains a citation, so the wrapper carries no bibliography and no bibliography entry.  No retained line contains a figure environment, an image-inclusion command or drawing code, so no image is delivered and the package declares no asset dependency.  Editorial formatting only, in addition: an AMS \texttt{amsart} preamble that restates the article's own declarations and macros used by the retained lines, namely the environments \texttt{assumption}, \texttt{theorem} on separate counters, together with the standard packages those lines need.  The retained environments print in this document as Assumption 1, Theorem 1 in order of appearance, whereas the article prints the same items with the numbering of its own full document.  The complete native source of the article as distributed in the arXiv e-print for this exact version is bundled unchanged as \texttt{sources/205508/source0.tex}.
\begin{align}\label{tso_match_condn}
 \frac{\sigma^2}{2(\eta_1 + \eta_2)} = \sigma_J^2 .
\end{align} 

\begin{assumption}\label{tso_assum_2}
   Let $f \in C^{4}(B_r(x_{\min}))$ for some $r > 0$ with $x_{\min}$ being the unique global minimizer. We assume $\nabla^2 f(x_{\min}) > 0$. We assume that there exist constants $\kappa > 0$ and  $\ell \geq 0 $ such that 
   \begin{align}\label{tso_eqn_coerc_assum}
       f(x) \geq \kappa |x|^2 - \ell, \quad  x \in \mathbb{R}^{d}.
   \end{align}
\end{assumption}

\begin{theorem}[Quantitative Laplace principle for the steady state]
\label{thm:steady_state_laplace}
Let the variance-matching condition \eqref{tso_match_condn} hold, such that the conditional steady-state law of the explorers is Gaussian with variance $\sigma_\star^2 = \sigma_J^2$, and let $c_0 = 1/(2\sigma_\star^2)$. Let the objective function $f$ satisfy Assumption~\ref{tso_assum_2}. Define the mapping $T_\alpha: \mathbb{R}^d \to \mathbb{R}^d$ as
\begin{align}
    T_\alpha(m) = \frac{\int_{\R^d} x e^{-\alpha f(x)} e^{-c_0 |x - m|^2} \der x}{\int_{\mathbb{R}^d} e^{-\alpha f(x)} e^{-c_0 |x - m|^2} \der x}.
\end{align}
Then, for all sufficiently large $\alpha > 0$, there exists a self-consistent steady-state consensus point $m_\alpha \in \mathbb{R}^d$ satisfying 
\begin{align}\label{tso_eqn_3.15}
    m_\alpha = T_\alpha(m_\alpha).
\end{align}
Moreover, there exists a constant $C > 0$, independent of $\alpha$, such that the distance between $m_\alpha$ and the global minimum satisfies the bound
\begin{align}
    |m_\alpha - x_{\min}| \le \frac{C}{\alpha}.
\end{align}
\end{theorem}

\begin{proof}
Without loss of generality, we may assume $f(x_{\min}) = 0$, as shifting $f$ by a constant cancels out in the definition of $T_\alpha(m)$.

Due to \eqref{tso_eqn_coerc_assum}, the mapping $T_\alpha$ is well-defined. Moreover, if we take a sequence $m_n $ converging to $m$ then $e^{-c_0 |x - m_n|^2} \to e^{-c_0 |x - m|^2}$. Then, again due to  \eqref{tso_eqn_coerc_assum}, dominated convergence theorem can be applied to obtain $T_\alpha (m_n) \to T_{\alpha}(m) $ as $n \to \infty$. Therefore, $T_\alpha :\mathbb{R}^{d} \to \mathbb{R}^d$ is a continuous mapping. 


We now prove that for sufficiently large $R$ and sufficiently large $\alpha$, the following holds:
\begin{align*}
 \big( ( m - x_{\min}) \cdot (T_\alpha(m) - x_{\min})\big) \leq |m - x_{\min}|^2, 
\end{align*}
whenever $|m - x_{\min}| \geq R$. Let $y = x - x_{\min}$ and $\varepsilon = m - x_{\min}$, then using this notation, we have
\begin{align*}
    T_{\alpha}(m) - x_{\min} &= \frac{\int_{\mathbb{R}^d} y e^{-\alpha f(x_{\min} + y)} e^{-c_0 |y - \varepsilon|^2} \der y}{\int_{\mathbb{R}^d} e^{-\alpha f(x_{\min} + y)} e^{-c_0 |y - \varepsilon|^2} \der y}, \\ 
   \text{and}\quad  \big( \varepsilon \cdot (T_{\alpha}(m) - x_{\min})\big) &= \frac{\int_{\mathbb{R}^d} (\varepsilon \cdot y) e^{-\alpha f(x_{\min} + y)} e^{-c_0 |y - \varepsilon|^2} \der y}{\int_{\mathbb{R}^d} e^{-\alpha f(x_{\min} + y)} e^{-c_0 |y - \varepsilon|^2} \der y}.
\end{align*}
Consider the integral 
\begin{align*}
    \int_{\mathbb{R}^d} (\varepsilon \cdot y) e^{-\alpha f(x_{\min} + y)} e^{-c_0 |y - \varepsilon|^2} \der y & = \int_{|y| \leq |\varepsilon|/2 } (\varepsilon \cdot y) e^{-\alpha f(x_{\min} + y)} e^{-c_0 |y - \varepsilon|^2} \der y \nonumber 
    \\   &   \quad + \int_{|y| > |\varepsilon|/2 } (\varepsilon \cdot y) e^{-\alpha f(x_{\min} + y)} e^{-c_0 |y - \varepsilon|^2} \der y.
\end{align*}
Therefore, we get
\begin{align}\label{tso_eqn_2.7}
     \big( \varepsilon \cdot (T_{\alpha}(m) - x_{\min})\big) \leq  \frac{|\varepsilon|^2}{2 } + \frac{|\varepsilon|}{Z_\alpha(\varepsilon)} \int_{|y| > |\varepsilon|/2 }  |y| e^{-\alpha f(x_{\min} + y)} e^{-c_0 |y - \varepsilon|^2} \der y,
\end{align}
where $Z_\alpha(\varepsilon) := \int_{\mathbb{R}^d} e^{-\alpha f(x_{\min} + y)} e^{-c_0 |y - \varepsilon|^2} \der y $. 


For subsequent analysis, we need a lower bound for the denominator $Z_\alpha(\varepsilon)$. Since $f \in C^4$ near $x_{\min}$, $\nabla f(x_{\min}) = 0$, and $\nabla^2 f(x_{\min}) > 0$, there are constants $\hat{r} > 0$ and $ \hat{C} > 0$ such that
\begin{equation}\label{tso_eqn_2.8}
f(x_{\min} + y) \leq \hat C |y|^2 
\end{equation}
whenever $|y| \leq \hat{r}$. For sufficiently large $\alpha$, $|y| \le \alpha^{-1/2}$ implies $|y| \le \hat r$, which in turn guarantees
\begin{equation}\label{tso_eqn_2.9}
\alpha f(x_{\min} + y) \leq \hat{C}. 
\end{equation}
On the ball $|y| \leq \alpha^{-1/2}$, we also have 
\begin{equation} \label{tso_eqn_2.10}
e^{-c_0|y-\varepsilon|^2} \geq \exp\left(-c_0|\varepsilon|^2 - 2c_0|\varepsilon|\alpha^{-1/2} - c_0\alpha^{-1}\right),
\end{equation}
since $|y - \varepsilon|^{2} \leq |\varepsilon|^{2} + 2|\varepsilon||y| + |y|^2 \leq |\varepsilon|^2 + 2|\varepsilon|\alpha^{-1/2} + \alpha^{-1}$. Therefore, using \eqref{tso_eqn_2.9} and \eqref{tso_eqn_2.10}, we obtain
\begin{align}\label{tso_eqn_2.11}
Z_\alpha(\varepsilon) \geq C\alpha^{-d/2} \exp\left(- c_0|\varepsilon|^2 - 2c_0|\varepsilon|\alpha^{-1/2} - c_0\alpha^{-1}\right),
\end{align}
where $C>0$ is independent of $\alpha$.  We now obtain a bound on the integral in the second term on right hand side of inequality \eqref{tso_eqn_2.7}. 
%\begin{align} \int_{|y| > |\varepsilon|/2} |y| e^{-\alpha f(x_* + y)} e^{-c_0|y - \varepsilon|^2} \, dy. \end{align}
Since $|x|^2 = |x_{\min} + y|^2 \geq \frac{1}{2}|y|^2 - |x_{\min}|^2$, we have $f(x_{\min} + y) \geq \kappa |x_{\min} + y|^2 - \ell \ge \frac{\kappa}{2}|y|^2 - \hat{\ell}$ where $\hat{\ell} = \kappa |x_{\min}|^2 + \ell$. 
Therefore,
\begin{align}\label{tso_eqn_2.12}
 \int_{|y| > |\varepsilon| /2} |y| e^{-\alpha f(x_{\min} + y)} e^{-c_0|y - \varepsilon|^2} \der y \le C(1 + |\varepsilon|) \exp\left(\alpha \hat{\ell} - \frac{\alpha \kappa}{8}|\varepsilon|^2\right), 
\end{align} 
since $ \int_{\mathbb{R}^d} |y| e^{-c_0|y - \varepsilon|^2} \der y \le C(1 + |\varepsilon|)$ and $e^{-\alpha f(x_{\min} + y)} \leq e^{\alpha \hat{\ell}} e^{-\alpha \frac{\kappa}{2}| y|^2 }$.
Combining \eqref{tso_eqn_2.12} with \eqref{tso_eqn_2.11} gives
\begin{align*} 
\frac{|\varepsilon|}{Z_\alpha(\varepsilon)}  \int_{|y| > |\varepsilon| /2} |y| e^{-\alpha f(x_{\min} + y)} e^{-c_0|y - \varepsilon|^2} \der y \leq C|\varepsilon|(1 + |\varepsilon|) \alpha^{d/2} \exp\left(\alpha \hat \ell - \big(\frac{\alpha \kappa}{8} - c_0\big)|\varepsilon|^2 + C|\varepsilon|\alpha^{-1/2}\right).
\end{align*}
Choose $R > 0$ large enough that $\frac{\kappa}{8}R^2 > 2\hat{\ell}$. This gives 
\begin{align*} 
\alpha \hat{\ell} - \frac{\alpha \kappa}{8}R^2 \leq -\frac{\alpha \kappa}{16}R^2. 
\end{align*}
Therefore, for $|\varepsilon| \geq R$, $C |\varepsilon|(1+|\varepsilon|)\alpha^{d/2} \exp\left(-\frac{\alpha \kappa}{16}|\varepsilon|^2 + c_0|\varepsilon|^2 + 2 c_0 |\varepsilon|\alpha^{-1/2}\right) \to 0 $
as $\alpha \to \infty$. Hence, with sufficiently large $\alpha$, we can choose $\theta \in (0, 1/2)$ such that 
\begin{align*} 
\frac{|\varepsilon|}{Z_\alpha(\varepsilon)} \int_{|y| > |\varepsilon| /2} |y| e^{-\alpha f(x_{\min} + y)} e^{-c_0|y - \varepsilon|^2} \der y \leq (1-2\theta)\frac{|\varepsilon|^2}{2}
\end{align*}
for $|\varepsilon| \geq R$. Therefore,
\begin{align*}
\big(\varepsilon \cdot (T_\alpha(m) - x_{\min}) \big)\leq \frac{|\varepsilon|^2}{2} + (1 -2 \theta)\frac{|\varepsilon|^2}{2} = (1 -\theta)|\varepsilon|^2.
\end{align*}
Since $\varepsilon = m - x_{\min}$, this proves that for sufficiently large $R$ and sufficiently large $\alpha$,
\begin{equation} \label{tso_eqn_2.17}
\big( (m - x_{\min} ) \cdot (T_\alpha(m) - x_{\min})\big) \leq (1-\theta)|m - x_{\min}|^2,
\end{equation}
whenever $ |m - x_{\min}| \geq R$. Here $\theta \in (0, 1/2)$. 

We now use this estimate to obtain a fixed point. Let
$P:\mathbb R^d\to \bar{B}_R(x_{\min})$ be the radial projection onto
the closed ball centered at $x_{\min}$:
\begin{equation*}
    P(z)
    =
    \begin{cases}
    z, & |z-x_{\min}|\le R,\\
    x_{\min}     +     R\dfrac{z-x_{\min}}{|z-x_{\min}|},
    & |z-x_{\min}|>R.
    \end{cases}
\end{equation*}
Then $ P\circ T_\alpha:
    \bar{B}_R(x_{\min})
    \to     \bar{B}_R(x_{\min}) $ is continuous. By Brouwer's fixed-point theorem, there exists
$m_\alpha\in \bar{B}_R(x_{\min})$ such that
\begin{align*}
    m_\alpha=P(T_\alpha(m_\alpha)).
\end{align*}
We show  that $m_\alpha$ is indeed a fixed point of $T_\alpha$.
First consider the case when $|T_\alpha(m_\alpha)-x_{\min}|\le R$. Then, we simply have 
\begin{equation*}
    m_\alpha = P(T_\alpha(m_\alpha)) =     T_\alpha(m_\alpha).
\end{equation*}
It remains to deal with the case when $|T_\alpha(m_\alpha)-x_{\min}|>R$. 
In this case, by the definition of $P$, we have
\begin{align*}
     m_\alpha
    =  x_{\min} + R\frac{
        T_\alpha(m_\alpha)-x_{\min}}{|T_\alpha(m_\alpha)-x_{\min}|}.
\end{align*}
Hence $|m_\alpha-x_{\min}|=R$ and $m_\alpha-x_{\min}$ points in the same direction as
$T_\alpha(m_\alpha)-x_{\min}$. Equivalently, $T_\alpha(m_\alpha)-x_{\min} = \hat{a}(m_\alpha-x_{\min})$
 with $\hat{a} = \frac{|T_\alpha(m_\alpha)-x_{\min}|}{R} >1$. 
Therefore
\begin{align*}
    \big((m_\alpha-x_{\min})\cdot
    (T_\alpha(m_\alpha)-x_{\min}) \big)  &=
    \hat a |m_\alpha-x_{\min}|^2 = \hat{a} R^2 > R^2.
\end{align*}
This contradicts the estimate already proved i.e. \eqref{tso_eqn_2.17} whenever $|m-x_{\min}|=R$.
Therefore, $|T_\alpha(m_\alpha)-x_{\min}|>R$ is impossible.  Hence 
\begin{equation*}
    |T_\alpha(m_\alpha)-x_{\min}|\leq R,
\end{equation*} and
\begin{align}
    m_\alpha=T_\alpha(m_\alpha).
\end{align}
Therefore, there exists a self-consistent solution $m_\alpha$ and $m_\alpha \in \bar{B}_{R}(x_{\min})$ due to strict inward-pointing estimate \eqref{tso_eqn_2.17} outside the ball.




We now prove that every self-consistent solution satisfies the sharp estimate
\begin{align}
    |m_\alpha-x_{\min}|\le \frac{C}{\alpha}
\end{align}
for all sufficiently large $\alpha$. Let $m_\alpha$ be any solution of the self-consistency equation $ m_\alpha=T_\alpha(m_\alpha)$.  Let $\delta_\alpha:=m_\alpha-x_{\min}$ and therefore
\begin{equation*}
    \delta_\alpha = \frac{
    \int_{\mathbb R^d} y e^{-\alpha f(x_{\min}+y)}e^{-c_0|y-\delta_\alpha|^2}\der y}{ \int_{\mathbb R^d}
    e^{-\alpha f(x_{\min}+y)}
    e^{-c_0|y-\delta_\alpha|^2}\der y}.
\end{equation*}
We now use the uniform boundedness obtained above to prove the sharp rate. Since $|\delta_\alpha|\leq R$, $b:=2c_0\delta_\alpha$
is bounded uniformly in $\alpha$, i.e.
$
    |b|\le 2c_0R.
$ We set $B := 2 c_0 R$. 
Using $|y-\delta_\alpha|^2 = |y|^2-2(y\cdot\delta_\alpha) +|\delta_\alpha|^2$, 
we get
\begin{align*}
    \delta_\alpha = \frac{ \int_{\mathbb R^d}
    y e^{-\alpha f(x_{\min}+y)}
    e^{(b\cdot y)}
    e^{-c_0|y|^2}
    \der y}{    \int_{\mathbb R^d}
    e^{-\alpha f(x_{\min}+y)}
    e^{(b\cdot y)}
    e^{-c_0|y|^2}
    \der y}.
\end{align*}
For $|b|\le 2c_0 R$, denote
\begin{align*}
    \mathrm{D}_\alpha(b)
   & :=
    \int_{\mathbb R^d}
    e^{-\alpha f(x_{\min}+y)}
    e^{(b\cdot y)}
    e^{-c_0|y|^2}
    \der y,\\ 
    \mathrm{N}_\alpha(b)
   & :=
    \int_{\mathbb R^d}
    y\,e^{-\alpha f(x_{\min}+y)}
    e^{(b\cdot y)}
    e^{-c_0|y|^2}
    \der y.
\end{align*}
We now prove
\begin{equation}
    \mathrm{D}_\alpha(b)\ge c\alpha^{-d/2}\quad \text{and}\quad
    |\mathrm{N}_\alpha(b)|\le C\alpha^{-(d+2)/2},
\end{equation}
uniformly for $|b|\le 2c_0 R$. If we choose sufficiently large $\alpha $ such that $\alpha^{-1/2} \leq \hat{r}$ , we ascertain (see \eqref{tso_eqn_2.8})
\begin{equation*}
    \alpha f(x_{\min}+y)\le \hat C,
\end{equation*}
whenever $|y| \leq \alpha^{-1/2}$. Moreover, using $|b|\le 2c_0 R$, we obtain
\begin{align}\label{tso_eqn_2.29}
    \mathrm{D}_\alpha(b)
    &\ge
    \int_{|y|\le \alpha^{-1/2}}
    e^{-\alpha f(x_{\min}+y)}
    e^{(b\cdot y)}
    e^{-c_0|y|^2} \der y     \geq     c\alpha^{-d/2}.
\end{align}


We next estimate $\mathrm{N}_\alpha(b)$. We first fix a constant $r>0$ sufficiently small. Under Assumption~\ref{tso_assum_2}, there exists $\eta>0$ such that
\begin{equation*}
    f(x_{\min}+y)\ge \eta,
    \qquad |y|\ge r.
\end{equation*}
Therefore, 
\begin{align}\label{tso_eqn_2.31}
    &\int_{|y|\ge r}
    |y|e^{-\alpha f(x_{\min}+y)}
    e^{(b\cdot y)}
    e^{-c_0|y|^2}
    \der y \leq
    e^{-\alpha\eta}
    \int_{\mathbb R^d}
    |y|e^{B|y|}e^{-c_0|y|^2}\der y
    \leq    Ce^{-\alpha\eta},
\end{align}
where $C$ is independent of $\alpha$. It remains to estimate the bound when  $|y|<r$.   Let $y = \frac{z}{\sqrt{\alpha}}$ therefore $\der y = \alpha^{-d/2} \der z $. 
We have 
\begin{align}
    \int_{|y|<r}
    y\,&e^{-\alpha f(x_{\min}+y)}
    e^{(b\cdot y)}
    e^{-c_0|y|^2}
    \der y     \nonumber \\  & =  \alpha^{-(d+1)/2}
    \int_{|z|<r\sqrt{\alpha}} z    \exp\left(
        -\alpha f\left(x_{\min}+\frac{z}{\sqrt{\alpha}}\right)+
        \frac{(b\cdot z)}{\sqrt{\alpha}} - \frac{c_0|z|^2}{\alpha}
    \right)\der z. 
    \end{align}
    Using Taylor's theorem, we get
\begin{align}
f(x_{\min}+y) = \frac{1}{2} y^\top H y + \frac{1}{6} D^3 f(x_{\min})[y,y,y] + \mathcal{O}(|y|^4),
\end{align}
where $
    H:=\nabla^2 f(x_{\min})>0$. Choose $0<\gamma<1/8$. Then, for $|z|\le \alpha^\gamma$, we have
\begin{align*}
    -\alpha f\left(x_{\min}+\frac z{\sqrt\alpha}\right) + \frac{b\cdot z}{\sqrt\alpha} -
    \frac{c_0|z|^2}{\alpha}
    &=  -\frac{1}{2} z^\top H z
    +  \frac{1}{\sqrt\alpha}A(z)
    + Q_\alpha(z),
\end{align*}
where $ A(z)
    := (b\cdot z) - \frac{1}{6} D^3f(x_{\min})[z,z,z]$, $|Q_\alpha(z)| \leq C\frac{|z|^2+|z|^4}{\alpha}$ and uniformly for $|b|\le B$, $|A(z)|\le C(1+|z|^3)$. Since $\gamma<1/8$, $\frac{1}{\sqrt\alpha}A(z) + Q_\alpha(z)$
tends to zero uniformly on $|z|\leq \alpha^\gamma$. Hence Taylor's theorem, uniformly for \(|b|\le B\), gives
\begin{equation}\label{tso_eqn_3.31}
    \exp\left(
        \frac1{\sqrt\alpha}A(z)+Q_\alpha(z)
    \right) = 1 + \frac{1}{\sqrt\alpha}A(z) +    \mathcal{O}\left(\frac{1+|z|^8}{\alpha}\right)
\end{equation}
on \(|z|\le \alpha^\gamma\). 
Since  $\int_{|z| < \alpha^\gamma} z e^{-\frac{1}{2}z^{\top}H z} \der z = 0$,  we get
\begin{equation}\label{tso_eqn_2.40}
   \bigg| \int_{|z|<\alpha^\gamma}
    ye^{-\alpha f(x_{\min}+y)}
    e^{(b\cdot y)}
    e^{-c_0|y|^2}
    \der y\bigg| \leq
    C\alpha^{-(d+2)/2},
\end{equation}
uniformly for $|b|\leq 2c_0 R$.  As in \eqref{tso_eqn_2.8}, for sufficiently small $r >0$, there is a constant $\tilde{C}$ such that
\begin{align*}
    f(x_{\min} + y) \geq \tilde{C} |y|^2,\qquad |y|\leq r.
\end{align*}
For all sufficiently large $\alpha$, since $|z|>\alpha^\gamma$, $\frac{B|z|}{\sqrt\alpha}
    \leq \frac{\tilde C}{2}|z|^2$. Therefore, for sufficiently large $\alpha$, we get
\begin{align}\label{eqn_tso_3.33}
        \alpha^{-(d+1)/2}
    \int_{\alpha^\gamma<|z|<r\sqrt\alpha}
    |z|\exp\left(-\tilde{C} |z|^2
        +
        \frac{B|z|}{\sqrt\alpha}
    \right)\der z & \leq  
    C\alpha^{-(d+1)/2}
    \int_{|z|>\alpha^\gamma}
    |z|e^{-\frac{\tilde{C}}{2}|z|^2}\,\der z  \nonumber 
    \\   &  \leq  C\alpha^{-(d+2)/2}.
\end{align}
Combining \eqref{tso_eqn_2.31}, \eqref{tso_eqn_2.40} and \eqref{eqn_tso_3.33}, we get
\begin{align}\label{tso_eqn_2.41}
    |\mathrm{N}_\alpha(b)| \leq
    C\alpha^{-(d+2)/2}.
\end{align}
Therefore, using \eqref{tso_eqn_2.29} and \eqref{tso_eqn_2.41}, we obtain
\begin{align}
 |m_\alpha - x_{\min}| =    \left|\frac{N_\alpha(b)}{D_\alpha(b)}\right|
    \leq
    \frac{
        C\alpha^{-(d+2)/2} }{        c\alpha^{-d/2}} =
    \frac{C}{\alpha}.
\end{align}
\end{proof}

\end{document}
